% =====================================================================
%  CHECKPOINT 04 -- the finished document
%
%  This is where the workshop ends up.
%  Keep it: it is a working skeleton for your next report, and almost
%  everything you will need for a term paper or a thesis is in here.
% =====================================================================
%  A NOTE ON WHAT THIS DOCUMENT IS
%  It is a worked example, not a real experiment. Nobody measured any
%  coffee: the numbers and the plot were invented for this workshop.
%  Your job today is to typeset the report, not to run it.
% =====================================================================
%  HOUSE STYLE USED IN THESE FILES
%    * one sentence per line -- a long one may break at a comma
%    * \label goes on the same line as the \section it belongs to
%    * four spaces per indent level
%  None of this changes the output by a single pixel. It changes how
%  readable the SOURCE is, and how useful `git diff` is later.
% =====================================================================

\documentclass[a4paper,11pt]{article}

% ---------------------------------------------------------------------
%  Preamble
% ---------------------------------------------------------------------
\usepackage[T1]{fontenc}        % proper 8-bit output encoding
\usepackage[english]{babel}     % hyphenation and language conventions
\usepackage{newtxtext}          % Times for the text
\usepackage{parskip}            % blank line between paragraphs, no indent
\usepackage{amsmath}            % align, cases, \text, decent spacing
\usepackage{newtxmath}          % ...and Times for the maths to match
\usepackage{booktabs}           % tables that do not look like spreadsheets
\usepackage{caption}            % sane space between a caption and its float
\usepackage{graphicx}           % \includegraphics

\usepackage[backend=biber,style=numeric]{biblatex}
\addbibresource{references.bib}

% Prefer a sans-serif document? Replace the two newtx lines above with:
%   \usepackage{helvet}
%   \renewcommand{\familydefault}{\sfdefault}
%
% Worth adding next, in this order:
%   \usepackage{siunitx}        % \qty{65}{\degreeCelsius} -- see below
%   \usepackage{microtype}      % quietly better justification
%   \usepackage{hyperref}       % clickable refs -- load this one last

\title{How Fast Does Your Coffee Go Cold?}
\author{Your Name}
\date{\today}

\begin{document}

\maketitle
\tableofcontents

% =====================================================================
\section{Introduction}\label{sec:introduction}

A freshly poured coffee is too hot to drink, and a forgotten one is too cold.
Somewhere in between there is a window of a few minutes in which the drink is actually pleasant,
and this report sets out to find it.
The underlying physics is \emph{Newton's law of cooling},
which says that a body loses heat in proportion to how much hotter it is than its surroundings.

Three \textbf{everyday vessels} are compared,
referred to throughout by the short labels \texttt{mug}, \texttt{glass} and \texttt{flask}.
The question is worth asking precisely because the model is so thin.
It has exactly one free parameter per vessel,
so a single afternoon is enough to test it properly.

This report sets out to:

\begin{itemize}
    \item state the cooling model and what its one parameter means,
    \item compare the cooling curves of the three vessels, and
    \item work out how long each of them keeps a drink in the pleasant range.
\end{itemize}

% =====================================================================
\section{Theory}\label{sec:theory}

The model is described by three quantities:

\begin{description}
    \item[$T_0$] temperature of the drink at the start
    \item[$T_\mathrm{env}$] temperature of the room
    \item[$\tau$] time constant of the vessel: the time it takes to close
        roughly two thirds of the gap to room temperature
\end{description}

Newton's law of cooling states that the temperature difference to the surroundings decays
exponentially \cite{halliday2018}, so that

\begin{equation}
    T(t) = T_\mathrm{env} + \left( T_0 - T_\mathrm{env} \right) \mathrm{e}^{-t/\tau} .
    \label{eq:cooling}
\end{equation}

Rearranging for the time gives the waiting time until the drink has reached a chosen
temperature $T$,

\begin{align}
    \frac{T - T_\mathrm{env}}{T_0 - T_\mathrm{env}}
        &= \mathrm{e}^{-t/\tau} \notag \\
    t &= -\tau \ln \left( \frac{T - T_\mathrm{env}}{T_0 - T_\mathrm{env}} \right) .
    \label{eq:waiting}
\end{align}

Equations~\eqref{eq:cooling} and~\eqref{eq:waiting} are all the theory this report needs:
the first says what the curve looks like, the second says when to pick the cup up.

What counts as pleasant is of course a matter of taste.
This report uses

\[
    \operatorname{verdict}(T) =
    \begin{cases}
        \text{too hot}   & \text{if} \quad T > 65\,^{\circ}\mathrm{C} \\
        \text{drinkable} & \text{if} \quad 50\,^{\circ}\mathrm{C} \leq T \leq 65\,^{\circ}\mathrm{C} \\
        \text{too cold}  & \text{if} \quad T < 50\,^{\circ}\mathrm{C}
    \end{cases}
\]

% =====================================================================
\section{Setup}\label{sec:setup}

Each vessel starts with 300 ml of water at 85 degrees Celsius,
in a closed room at an ambient temperature of 21 degrees Celsius.
The three vessels are listed in Table~\ref{tab:vessels}.

A thermometer sits in the centre of the liquid, clear of the walls,
and the temperature is taken every 30 seconds for one hour.

The procedure is as follows:

\begin{enumerate}
    \item Heat 300 ml of water to 85 degrees Celsius.
    \item Pour it into the vessel and start the clock.
    \item Suspend the thermometer in the centre of the liquid.
    \item Record the temperature every 30 seconds for one hour.
\end{enumerate}

\begin{table}[htbp]
    \centering
    \caption{The three vessels and their time constants.}
    \label{tab:vessels}
    \begin{tabular}{llllr}
        \toprule
        Label             & Vessel     & Material       & Lid & $\tau$   \\
        \midrule
        \texttt{mug}      & Mug        & glazed ceramic & no  & 24 min   \\
        \texttt{glass}    & Glass      & borosilicate   & no  & 18 min   \\
        \texttt{flask}    & Travel mug & steel, vacuum  & yes & 310 min  \\
        \bottomrule
    \end{tabular}
\end{table}

% =====================================================================
\section{Results}\label{sec:results}

\begin{figure}[htbp]
    \centering
    \includegraphics[width=0.8\textwidth]{figures/cooling-curves}
    \caption{Cooling curves of the three vessels, with the model of
        Equation~\eqref{eq:cooling} drawn through them.
        The shaded band marks the drinkable range.}
    \label{fig:cooling}
\end{figure}

Figure~\ref{fig:cooling} shows the temperatures together with the model.
All three vessels follow the exponential closely;
what separates them is only the time constant listed in Table~\ref{tab:vessels}.

Feeding those time constants into Equation~\eqref{eq:waiting} gives the drinkable window directly.
The \texttt{mug} enters it about 9 minutes after pouring and leaves it after 19,
which is the ten-minute window most people will recognise.
The \texttt{glass} gives only seven minutes.
The \texttt{flask} is still too hot to drink almost two hours later,
which is either its great virtue or its great flaw,
depending on what you wanted from it \cite{incropera2017}.

% =====================================================================
\section{Conclusion}\label{sec:conclusion}

Newton's law of cooling describes all three vessels well over a full hour,
with a single time constant each.
The practical answer is that an open mug of coffee is worth drinking between roughly 9 and 19
minutes after it is poured,
and that insulation does not so much keep the drink pleasant as postpone the moment it becomes so.

\printbibliography

\end{document}
